\section{Οργανωσιακό Οικοσύστημα και Όρια των Αποθετηρίων}
\label{sec:appendix_ecosystem_scope}

Η υλοποίηση της διατριβής δεν αναπτύχθηκε σε ένα μονολιθικό αποθετήριο, αλλά σε ένα μικρό οργανωσιακό οικοσύστημα
γύρω από το \ENG{GitHub Organization} \texttt{\ENG{Evangelos-Komnis-Thesis-ICT-IHU-2025}}.
Το οποίο φιλοξενεί τον \ENG{SCORM Engine}, τον \ENG{player},
το \ENG{PHP SDK}, το \ENG{Laravel LMS} για λόγους δοκιμών, την κεντρική \ENG{docker} υποδομή, το αποθετήριο του παρόντος κειμένου,
και επιπλέον βοηθητικά αποθετήρια για \ENG{workflow templating}, \ENG{AI-assisted prompts} και παραδείγματα μαθημάτων.
Η εικόνα αυτή έχει σημασία επειδή αποκαλύπτει ότι η συγκεκριμένη υλοποίηση δεν στηρίζεται μόνο σε κώδικα, αλλά και σε ένα 
ρητά οργανωμένο περιβάλλον αυτοματοποίησης και συγγραφικής υποστήριξης.

Η βασική οργανωτική διάκριση είναι επομένως τριμερής. Πρώτον, υπάρχει το \ENG{runtime and integration perimeter}, όπου ανήκουν τα αποθετήρια
που σχετίζονται με την εκτέλεση και την ενσωμάτωση της πλατφόρμας. Δεύτερον, υπάρχει το \ENG{operational perimeter}, όπου ανήκουν η κεντρική
υποδομή, τα \ENG{workflow templates} και τα διάφορα \ENG{scripts} επαλήθευσης. Τρίτον, υπάρχει το \ENG{research and writing perimeter},
όπου εντάσσονται το αποθετήριο της διπλωματικής και τα \ENG{AI-assisted} συγγραφικά \ENG{artifacts}.
Το Σχήμα~\ref{fig:appendix_org_repository_map} συνοψίζει αυτή τη διάκριση, ενώ ο Πίνακας~\ref{tab:appendix_repository_evidence_matrix} 
αποδίδει για κάθε βασικό αποθετήριο τον τεκμηριωμένο ρόλο του και τη θέση του στο σύνολο.

\begin{figure}[H]
  \centering
\begin{chapterfigurebox}
  \begin{tikzpicture}[
      font=\footnotesize,
      node distance=1.7cm and 2.3cm,
      box/.style={draw, rounded corners, fill=gray!5, align=center, text width=3.5cm, minimum height=1.05cm},
      emph/.style={draw, rounded corners, fill=gray!12, align=center, text width=3.6cm, minimum height=1.05cm, font=\footnotesize\bfseries},
      boundary/.style={draw, dashed, rounded corners, inner sep=11pt},
      arrow/.style={-Latex, thick},
      support/.style={-Latex, thick, dashed}
    ]
    \node[box] (template) {\texttt{\ENG{github-actions-template}}\\\ENG{workflow generation}};
    \node[box, below=0.9cm of template] (ai) {\texttt{\ENG{ai-agents-instructions}}\\\ENG{AI support prompts}};
    \node[box, below=0.9cm of ai] (examples) {\texttt{\ENG{example-scorm-files}}\\\ENG{course corpus}};

    \node[emph, right=3.0cm of template] (engine) {\texttt{\ENG{scorm-engine}}};
    \node[emph, below=0.9cm of engine] (player) {\texttt{\ENG{player}}};
    \node[emph, right=of engine] (sdk) {\texttt{\ENG{scorm-engine-php-sdk}}};
    \node[emph, below=0.9cm of sdk] (lms) {\texttt{\ENG{example-lms-client}}};
    \node[emph, right=of player] (infra) {\texttt{\ENG{central-docker-infrastructure}}};

    \node[box, below=2.0cm of player] (thesis) {\texttt{\ENG{thesis-document-latex}}\\\ENG{research writing}};
    \node[box, below=2.0cm of infra] (orgmeta) {\texttt{\ENG{.github}}\\\ENG{organization metadata}};

    \draw[arrow] (lms) -- (sdk);
    \draw[arrow] (sdk) -- (engine);
    \draw[arrow] (player) -- (engine);
    \draw[arrow] (infra) -- (engine);
    \draw[arrow] (infra) -- (player);
    \draw[arrow] (infra) -- (lms);

    \draw[support] (template) -- (sdk);
    \draw[support] (template) -- (lms);
    \draw[support] (template) -- (engine);
    \draw[support] (template) -- (player);
    \draw[support] (ai) -- (thesis);
    \draw[support] (ai) -- (engine);
    \draw[support] (examples) -- (thesis);

    \node[boundary, fit=(engine)(player)(sdk)(lms)(infra)] (core) {};
    \node[font=\footnotesize, align=center, text width=5.6cm, above=0.16cm of core] {\ENG{core runtime and integration perimeter}};

    \node[boundary, fit=(template)(ai)(examples)(orgmeta)] (supporting) {};
    \node[font=\footnotesize, align=center, text width=6.2cm, above=0.16cm of supporting] {\ENG{supporting repositories and organizational assets}};
  \end{tikzpicture}
\end{chapterfigurebox}
  \caption{Οργανωσιακός χάρτης του τεχνικού οικοσυστήματος της διατριβής, με διάκριση ανάμεσα στον πυρήνα εκτέλεσης και στα υποστηρικτικά αποθετήρια.}
  \label{fig:appendix_org_repository_map}
\end{figure}


\begin{table}[H]
\centering
\caption{Χαρτογράφηση αποθετηρίων, τεκμηρίων και ρόλων στο συνολικό τεχνικό οικοσύστημα}
\label{tab:appendix_repository_evidence_matrix}
\begingroup
\footnotesize
\setlength{\tabcolsep}{3pt}
\renewcommand{\arraystretch}{1.20}
\begin{adjustbox}{center,max width=\textwidth}
\begin{tabular}{|p{3.0cm}|p{3.9cm}|p{4.8cm}|p{3.5cm}|}
\hline
Αποθετήριο & Κύριος ρόλος & Ισχυρότερα τεκμήρια & Ρόλος στα παραρτήματα \\
\hline
\texttt{\ENG{scorm-engine}} & Κεντρικός \ENG{backend} πυρήνας και επιφάνεια διεπαφών & \texttt{\ENG{pom.xml}}, \texttt{\ENG{application.yml}}, \texttt{\ENG{OpenApiConfig.java}}, \texttt{\ENG{V1\_\_init\_schema.sql}} & \ENG{API support}, \ENG{QA}, \ENG{coverage}, \ENG{release}, \ENG{observability} \\
\hline
\texttt{\ENG{player}} & \ENG{Browser-facing} υπηρεσία εκκίνησης και διανομής περιεχομένου & \texttt{\ENG{Dockerfile}}, \texttt{\ENG{package.json}}, \texttt{\ENG{LaunchPageRenderer.ts}}, \texttt{\ENG{ContentCacheService.ts}} & \ENG{runtime support tooling}, \ENG{coverage}, \ENG{containerization} \\
\hline
\texttt{\ENG{scorm-engine-php-sdk}} & Τυποποιημένο \ENG{client contract layer} προς το \ENG{engine} & \texttt{\ENG{composer.json}}, \texttt{\ENG{phpunit.xml}}, \texttt{\ENG{ApiHttpClient.php}}, \texttt{\ENG{RetryTransport.php}} & \ENG{contract support}, \ENG{QA}, \ENG{versioning} \\
\hline
\texttt{\ENG{example-lms-client}} & \ENG{Laravel LMS} αναφοράς και \ENG{integration consumer} & \texttt{\ENG{ScormEngineServiceProvider.php}}, \texttt{\ENG{ScormEngineGateway.php}}, \texttt{\ENG{LearnerController.php}} & \ENG{integration practice}, \ENG{CI}, περιορισμένη \ENG{test evidence} \\
\hline
\texttt{\ENG{central-docker-infrastructure}} & \ENG{Compose} υποδομή, \ENG{stores}, \ENG{smoke checks}, \ENG{logging pipeline} & \texttt{\ENG{docker-compose.yml}}, \texttt{\ENG{entrypoint.sh}}, \texttt{\ENG{smoke-check.sh}}, \texttt{\ENG{logstash.conf}} & \ENG{infrastructure}, \ENG{DevOps}, \ENG{reproducibility} \\
\hline
\texttt{\ENG{github-actions-template}} & Τοπικό \ENG{template} παραγωγής \ENG{workflow} και \ENG{repository metadata} & \texttt{\ENG{template\_repo\_agent.py}}, \texttt{\ENG{README.md}} & \ENG{engineering support tooling} \\
\hline
\texttt{\ENG{ai-agents-instructions}} & \ENG{AI-assisted prompts} για ανάπτυξη και συγγραφή & \texttt{\ENG{ai\_unit\_test\_writer\_agent\_prompt.md}}, \texttt{\ENG{ai\_openapi\_documentation\_writer\_prompt.md}} & \ENG{AI support workflow evidence} \\
\hline
\texttt{\ENG{thesis-document-latex}} & Συγγραφική υποδομή της διατριβής & \texttt{\ENG{thesis.tex}}, \ENG{chapter folders}, \texttt{\ENG{build.cmd}} & Μορφολογική και τεκμηριωτική υποστήριξη, όχι \ENG{runtime subsystem} \\
\hline
\end{tabular}
\end{adjustbox}
\endgroup

\end{table}
